% Options for packages loaded elsewhere
\PassOptionsToPackage{unicode}{hyperref}
\PassOptionsToPackage{hyphens}{url}
\PassOptionsToPackage{space}{xeCJK}
\documentclass[
  a4paper,11pt]{ltjsarticle}
\usepackage{xcolor}
\usepackage[margin=1in]{geometry}
\usepackage{amsmath,amssymb}
\usepackage{graphicx}
\setcounter{secnumdepth}{-\maxdimen} % remove section numbering
\usepackage{iftex}
\ifPDFTeX
  \usepackage[T1]{fontenc}
  \usepackage[utf8]{inputenc}
  \usepackage{textcomp}
\else
  \usepackage{unicode-math}
  \defaultfontfeatures{Scale=MatchLowercase}
  \defaultfontfeatures[\rmfamily]{Ligatures=TeX,Scale=1}
\fi
\usepackage{lmodern}
\ifPDFTeX\else
  \setmainfont[]{Times New Roman}
  \ifXeTeX
    \usepackage{xeCJK}
    \setCJKmainfont[]{Hiragino Mincho ProN}
  \fi
  \ifLuaTeX
    \usepackage[]{luatexja-fontspec}
    \setmainjfont[]{Hiragino Mincho ProN}
  \fi
\fi
\IfFileExists{upquote.sty}{\usepackage{upquote}}{}
\IfFileExists{microtype.sty}{%
  \usepackage[]{microtype}
  \UseMicrotypeSet[protrusion]{basicmath}
}{}
\makeatletter
\@ifundefined{KOMAClassName}{%
  \IfFileExists{parskip.sty}{%
    \usepackage{parskip}
  }{%
    \setlength{\parindent}{0pt}
    \setlength{\parskip}{6pt plus 2pt minus 1pt}}
}{% if KOMA class
  \KOMAoptions{parskip=half}}
\makeatother
\usepackage{longtable,booktabs,array}
\newcounter{none}
\usepackage{calc}
\usepackage{etoolbox}
\makeatletter
\patchcmd\longtable{\par}{\if@noskipsec\mbox{}\fi\par}{}{}
\makeatother
\IfFileExists{footnotehyper.sty}{\usepackage{footnotehyper}}{\usepackage{footnote}}
\makesavenoteenv{longtable}
\usepackage{bookmark}
\IfFileExists{xurl.sty}{\usepackage{xurl}}{}
\urlstyle{same}
\usepackage{hyperref}
\hypersetup{
  hidelinks,
  pdfcreator={LaTeX via pandoc}}
\setlength{\emergencystretch}{3em}
\providecommand{\tightlist}{%
  \setlength{\itemsep}{0pt}\setlength{\parskip}{0pt}}
\usepackage{amsthm}
\newtheorem{theorem}{定理}[section]
\newtheorem{proposition}[theorem]{命題}
\newtheorem{lemma}[theorem]{補題}
\newtheorem{corollary}[theorem]{系}
\theoremstyle{definition}
\newtheorem{definition}[theorem]{定義}
\theoremstyle{remark}
\newtheorem{remark}[theorem]{注意}
\newtheorem{observation}[theorem]{観察}
\begin{document}
\section{質量構造の考察------軸スケール値と符号付き面積による質量分析の枠組み}\label{ux8ceaux91cfux69cbux9020ux306eux8003ux5bdfux8ef8ux30b9ux30b1ux30fcux30ebux5024ux3068ux7b26ux53f7ux4ed8ux304dux9762ux7a4dux306bux3088ux308bux8ceaux91cfux5206ux6790ux306eux67a0ux7d44ux307f}

\subsection{------ {[}4{]} §9.11 の質量問題に対する予備的考察
------}\label{ux306eux8ceaux91cfux554fux984cux306bux5bfeux3059ux308bux4e88ux5099ux7684ux8003ux5bdf}

\textbf{著者}: 木原 範昭 (WF System Co., Ltd.)

\textbf{日付}: 2026年4月

\begin{center}\rule{0.5\linewidth}{0.5pt}\end{center}

\subsection{§0
本稿の位置づけ}\label{ux672cux7a3fux306eux4f4dux7f6eux3065ux3051}

{[}4{]} §9.11
において、質量比の定量的導出は未解決課題とされた。本稿は、この課題に対する本格的な分析に先立ち、{[}4{]}
の集合構造と補講1の符号付き面積 \(M_F\)
を用いた質量分析がどのような形をとりうるかを考察する。

本稿は予備的考察であり、具体的な質量公式の導出や質量値の予言を主張しない。考察の過程で得られた構造的観察と、最も単純な質量公式が定量的に不十分であることの確認を記録する。

本稿で考察する質量構造は、{[}4{]} 定義9.4 の静止質量写像
\(f_m^{\text{rest}}: (x, y, z, t, R, Q) \to \mathbb{R}_{\geq 0}\)
および定義9.5 の曲率エネルギー写像
\(f_m^{\text{curv}}: (x, y, z, t, R, Q, R_1) \to \mathbb{R}_{\geq 0}\)
の関数形を制約する試みとして位置づけられる。§2--§3 はフェルミオンの
\(f_m^{\text{rest}}\) の候補形式を検討し、§5.1 はゲージボソンの
\(f_m^{\text{curv}}\) との関連に触れる。

\textbf{参照論文}: - {[}3{]}
離散空間における中心投影の全球被覆と5次元背景空間の整数論的必然性（W5）
- {[}4{]} 6次元超直方体の集合構造とその組合せ論的性質（W8） - {[}5{]}
形不変波の相互作用（W11） - {[}S1{]}
符号付き面積の定式化------電荷構造の導出とスピン2配置の帰結（補講1）

\begin{center}\rule{0.5\linewidth}{0.5pt}\end{center}

\subsection{§1
軸スケール値の導入}\label{ux8ef8ux30b9ux30b1ux30fcux30ebux5024ux306eux5c0eux5165}

\subsubsection{1.1
非零軸の座標値}\label{ux975eux96f6ux8ef8ux306eux5ea7ux6a19ux5024}

{[}4{]} §1.1 および {[}3{]}
により、各座標軸は非零の値（実数値・整数値のいずれでもよい）をとりうる。{[}S1{]}
定義1.2 により、非零軸の値 \(v_i\) は符号付き面積 \(M = v_i \cdot v_j\)
を通じて物理的意味を持つ。

ヒッグスボソン（定常波、\(\kappa = 1\)）が \(xyz\) 3軸のうち \(z\)
軸を選択することにより（{[}4{]} §9.11.7）、\(\mathrm{SO}(3)\) 対称性が
\(\mathrm{SO}(2)\)
に自発的に破れる。この結果、3つの空間軸のスケール値に階層が生じる:

\[v_z \gg v_x \approx v_y\]

\(v_z\) はヒッグスの主振動軸であり直接結合を受ける。\(v_x, v_y\)
は間接結合のみを受けるため相対的に小さい。

\subsubsection{1.2
4種のスケール値}\label{ux7a2eux306eux30b9ux30b1ux30fcux30ebux5024}

{[}4{]} §9.11.1 により、\(xyzt\) 4軸の非零本数 \(k\)
に応じて4種のスケール値が区別される:

{\def\LTcaptype{none} % do not increment counter
\begin{longtable}[]{@{}clll@{}}
\toprule\noalign{}
\(k\) & 型 & 非零軸 & スケール \\
\midrule\noalign{}
\endhead
\bottomrule\noalign{}
\endlastfoot
1 & ヒッグス & \(z\) のみ & \(v_z\) \\
2 & フェルミオン & 2軸の組合せ & \(v_i \cdot v_j\) \\
3 & スピン1ボソン（一部） & 3軸 & \(v_i \cdot v_j \cdot v_k\) \\
4 & スピン2ボソン & 4軸 & \(v_x \cdot v_y \cdot v_z \cdot v_t\) \\
\end{longtable}
}

本稿では \(k = 2\)（フェルミオン）の質量構造に焦点を当てる。

\begin{center}\rule{0.5\linewidth}{0.5pt}\end{center}

\subsection{§2
フェルミオン質量と軸スケール値}\label{ux30d5ux30a7ux30ebux30dfux30aaux30f3ux8ceaux91cfux3068ux8ef8ux30b9ux30b1ux30fcux30ebux5024}

\subsubsection{\texorpdfstring{2.1 （空間,
\(t\)）型フェルミオンの質量}{2.1 （空間, t）型フェルミオンの質量}}\label{ux7a7aux9593-tux578bux30d5ux30a7ux30ebux30dfux30aaux30f3ux306eux8ceaux91cf}

荷電レプトンおよびクォークは（空間1軸, \(t\)）型の2-面を持つ（{[}4{]}
§9.9 表A）。最も単純な仮定として、質量が符号付き面積の絶対値
\(|M_F| = |v_{\text{空間}} \cdot v_t|\) に比例すると考える:

\[m \propto |v_{\text{空間}} \cdot v_t|\]

これは {[}4{]} 定義9.4 の \(f_m^{\text{rest}}\)
に対する最も単純な候補形式であり、フェルミオン面の2軸 \(i, j\) を用いて:

\[f_m^{\text{rest}}(\sigma) = C \cdot |v_i \cdot v_j| \qquad (\sigma \text{ のフェルミオン面が } (i, j) \text{ 型})\]

と表される。以下ではこの候補の帰結を検討する。

荷電レプトン（\(Q = 0\)）の場合:

{\def\LTcaptype{none} % do not increment counter
\begin{longtable}[]{@{}cccl@{}}
\toprule\noalign{}
世代 & 空間軸 & 質量 & 質量の構造 \\
\midrule\noalign{}
\endhead
\bottomrule\noalign{}
\endlastfoot
1 & \(x\) & \(m_e = 0.511\) MeV & \(\propto v_x \cdot v_t\) \\
2 & \(y\) & \(m_\mu = 105.66\) MeV & \(\propto v_y \cdot v_t\) \\
3 & \(z\) & \(m_\tau = 1{,}776.86\) MeV & \(\propto v_z \cdot v_t\) \\
\end{longtable}
}

\(v_t\)
は全世代で共通であるから、世代間の質量比は空間軸スケール比に等しい:

\[\frac{m_\tau}{m_e} = \frac{v_z}{v_x} \approx 3{,}477, \qquad \frac{m_\tau}{m_\mu} = \frac{v_z}{v_y} \approx 16.82, \qquad \frac{m_\mu}{m_e} = \frac{v_y}{v_x} \approx 206.8\]

\subsubsection{2.2 （空間,
空間）型フェルミオンの質量}\label{ux7a7aux9593-ux7a7aux9593ux578bux30d5ux30a7ux30ebux30dfux30aaux30f3ux306eux8ceaux91cf}

ニュートリノは（空間, 空間）型の2-面を持つ（{[}4{]} §9.9
表A）。同じ仮定のもとで:

{\def\LTcaptype{none} % do not increment counter
\begin{longtable}[]{@{}ccl@{}}
\toprule\noalign{}
フレーバー & 非零軸 & 質量の構造 \\
\midrule\noalign{}
\endhead
\bottomrule\noalign{}
\endlastfoot
\(\nu_e\) & \((x, y)\) & \(\propto v_x \cdot v_y\) \\
\(\nu_\mu\) & \((y, z)\) & \(\propto v_y \cdot v_z\) \\
\(\nu_\tau\) & \((x, z)\) & \(\propto v_x \cdot v_z\) \\
\end{longtable}
}

\subsubsection{2.3 構造的観察}\label{ux69cbux9020ux7684ux89b3ux5bdf}

\textbf{観察 2.1}. 荷電レプトンとニュートリノは、同一の3つのスケール値
\(v_x, v_y, v_z\) の異なる組合せとして表現される。荷電レプトンは
\((v_i, v_t)\) の積、ニュートリノは \((v_i, v_j)\) の積である。

\textbf{観察 2.2}. §2.1 の荷電レプトン質量から空間軸スケール比
\(v_x : v_y : v_z\)
が確定すると、ニュートリノのフレーバー間質量比はパラメータなしで定まる:

\[\frac{m_{\nu_\mu}}{m_{\nu_e}} = \frac{v_y \cdot v_z}{v_x \cdot v_y} = \frac{v_z}{v_x} = \frac{m_\tau}{m_e} \approx 3{,}477\]

\[\frac{m_{\nu_\tau}}{m_{\nu_e}} = \frac{v_x \cdot v_z}{v_x \cdot v_y} = \frac{v_z}{v_y} = \frac{m_\tau}{m_\mu} \approx 16.82\]

\[\frac{m_{\nu_\mu}}{m_{\nu_\tau}} = \frac{v_y \cdot v_z}{v_x \cdot v_z} = \frac{v_y}{v_x} = \frac{m_\mu}{m_e} \approx 206.8\]

\textbf{注記}. ここでの \(m_{\nu_e}, m_{\nu_\mu}, m_{\nu_\tau}\)
はフレーバー固有状態に対応する幾何学的質量であり、実験で測定される質量固有状態
\(m_1, m_2, m_3\) とはPMNS行列で関係づけられる。PMNS行列自体も同じ
\(v_x, v_y, v_z\)
の構造から制約を受ける（補講2参照）。実測のニュートリノ質量二乗差（\(\Delta m^2_{21} \approx 7.5 \times 10^{-5}\)
eV²、\(|\Delta m^2_{32}| \approx 2.5 \times 10^{-3}\)
eV²）と本モデルのフレーバー間比の関係は、PMNS行列の具体的な形に依存するため、本稿の範囲を超える。

\begin{center}\rule{0.5\linewidth}{0.5pt}\end{center}

\subsection{§3
Q軸の質量への寄与}\label{qux8ef8ux306eux8ceaux91cfux3078ux306eux5bc4ux4e0e}

\subsubsection{3.1
Q軸結合の導入}\label{qux8ef8ux7d50ux5408ux306eux5c0eux5165}

クォーク（\(Q \in \{1, 2, 4\}\)）の質量は、荷電レプトン（\(Q = 0\)）と同じ空間軸を占有しながら異なる質量を持つ。最も単純な仮定として、\(Q\)
軸が質量に乗法的な因子 \(f(Q)\) を与えるとする:

\[m_{\text{quark}} = f(Q) \cdot v_{\text{空間}} \cdot v_t\]

これは \(f_m^{\text{rest}}\) が \(Q\)
軸の値に依存する場合の候補形式であり、§2.1 の
\(f_m^{\text{rest}}(\sigma) = C \cdot |v_i \cdot v_j|\) を
\(f_m^{\text{rest}}(\sigma) = f(Q) \cdot |v_i \cdot v_j|\)
に一般化したものに相当する。\(f(Q)\)
が全世代・全型で単一の定数であれば、クォークと対応する荷電レプトンの質量比は世代に依存しないはずである。

\subsubsection{3.2
実験値との比較}\label{ux5b9fux9a13ux5024ux3068ux306eux6bd4ux8f03}

{\def\LTcaptype{none} % do not increment counter
\begin{longtable}[]{@{}ccc@{}}
\toprule\noalign{}
世代 & 上型: \(m_q / m_\ell\) & 下型: \(m_q / m_\ell\) \\
\midrule\noalign{}
\endhead
\bottomrule\noalign{}
\endlastfoot
3 & \(m_t / m_\tau = 97.2\) & \(m_b / m_\tau = 2.35\) \\
2 & \(m_c / m_\mu = 12.0\) & \(m_s / m_\mu = 0.884\) \\
1 & \(m_u / m_e = 4.23\) & \(m_d / m_e = 9.14\) \\
\end{longtable}
}

\textbf{観察 3.1}. \(f(Q)\)
は定数ではない。世代間で大きく変動し、上型と下型でも異なる。

\textbf{観察 3.2}. \(f(Q)\) の世代依存性には規則性がある:
上型クォークでは世代が上がるほど \(f\)
が急増し（\(4.23 \to 12.0 \to 97.2\)）、下型クォークでは非単調（\(9.14 \to 0.884 \to 2.35\)）である。

\textbf{観察 3.3}. このパターンは、\(f(Q)\) が \(Q\)
軸の値のみに依存するのではなく、空間軸スケール \(v_i\)
にも依存すること------すなわち
\(f = f(Q, v_i)\)------を示唆する。質量公式
\(m = v_i \cdot v_t \cdot f(Q)\) の \(f\)
が定数であるという最も単純な仮定は不十分であり、\(v_i\)
との結合の関数形が質量問題の核心である。

\subsubsection{3.3
セクター内世代間質量比の比較}\label{ux30bbux30afux30bfux30fcux5185ux4e16ux4ee3ux9593ux8ceaux91cfux6bd4ux306eux6bd4ux8f03}

§5.2 予測1 の検証として、各セクター内の世代間質量比を比較する:

{\def\LTcaptype{none} % do not increment counter
\begin{longtable}[]{@{}lccc@{}}
\toprule\noalign{}
質量比 & 荷電レプトン & 上型クォーク & 下型クォーク \\
\midrule\noalign{}
\endhead
\bottomrule\noalign{}
\endlastfoot
\(m_3 / m_1\) & \(3{,}477\) & \(79{,}981\) & \(895\) \\
\(m_2 / m_1\) & \(206.8\) & \(588\) & \(20.0\) \\
\(m_3 / m_2\) & \(16.82\) & \(136\) & \(44.8\) \\
\end{longtable}
}

\textbf{観察 3.4}. \(f(Q)\)
が定数ならば各行の3列は同一値となるはずであるが、実測はすべて異なる。特に
\(m_3/m_1\) が荷電レプトンで \(3{,}477\) に対し上型クォークでは
\(79{,}981\) であることは、\(f(Q, v_i)\) が \(v_i\)
に対して強い非線形依存性を持つことを示している。

\begin{center}\rule{0.5\linewidth}{0.5pt}\end{center}

\subsection{§4
任意パラメータの数え上げ}\label{ux4efbux610fux30d1ux30e9ux30e1ux30fcux30bfux306eux6570ux3048ux4e0aux3052}

\subsubsection{4.1
標準模型の任意パラメータ}\label{ux6a19ux6e96ux6a21ux578bux306eux4efbux610fux30d1ux30e9ux30e1ux30fcux30bf}

{[}4{]}
の論文4（W4）により、標準模型は19個の任意パラメータを含む。このうち質量に関連するものは:

\begin{itemize}
\tightlist
\item
  フェルミオン質量: 9個（\(e, \mu, \tau, u, c, t, d, s, b\)）
\item
  ニュートリノ質量: 実測未確定（少なくとも2個の質量二乗差）
\item
  CKM行列: 4個（3角度 + 1位相）
\item
  PMNS行列: 4個（3角度 + 1位相）
\item
  ヒッグス質量: 1個
\end{itemize}

\subsubsection{4.2
本枠組みでのパラメータ構造}\label{ux672cux67a0ux7d44ux307fux3067ux306eux30d1ux30e9ux30e1ux30fcux30bfux69cbux9020}

本稿の考察が示唆するパラメータ構造は以下の通りである:

\textbf{基本パラメータ}: - \(v_x, v_y, v_z\):
空間軸スケール（3個。荷電レプトン質量から決定可能） - \(v_t\):
時間軸スケール（1個。ヒッグス質量との関係から決定可能性あり）

\textbf{導出される量}: - 荷電レプトン質量比: \(v_x : v_y : v_z\)
から決定（追加パラメータなし） - ニュートリノ質量比: 同じ
\(v_x : v_y : v_z\) から決定（追加パラメータなし） - PMNS混合角:
\(v_x, v_y, v_z\) とヒッグスの軸選択から定性的に決定（補講2） -
CKM混合角: 同上

\textbf{未決定の量}: - \(f(Q, v_i)\):
Q軸と空間軸の結合関数の関数形（§3.2 により定数ではない）。これは
\(f_m^{\text{rest}}\)（{[}4{]}
定義9.4）の具体的な関数形を決定する問題に相当する

\textbf{観察 4.1}.
質量に関連する19個以上のパラメータのうち、本枠組みで独立に残るのは: -
空間軸スケール比2個（\(v_y/v_x\), \(v_z/v_x\)）+ 全体スケール1個 -
Q軸結合関数 \(f(Q, v_i)\) の関数形

\(f(Q, v_i)\)
の関数形が確定すれば、自由パラメータは大幅に削減される可能性がある。\(f\)
が定数であれば1個まで削減されるが、§3.2
の実験値はこの最も単純な場合を否定している。

\begin{center}\rule{0.5\linewidth}{0.5pt}\end{center}

\subsection{§5 考察}\label{ux8003ux5bdf}

\subsubsection{5.1
質量公式の候補}\label{ux8ceaux91cfux516cux5f0fux306eux5019ux88dc}

§3 の結果は、正しい質量公式が \(m = v_i \cdot v_j\)
の単純な積ではないことを示している。{[}5{]} §8
では、質量の幾何学的起源として定常波（ヒッグス）による伝搬速度比
\(|k_x / k_t|\) の修正（\(\delta k_H\)）が示されている。\(\delta k_H\)
がフェルミオンの空間軸とヒッグスのz軸の\textbf{結合強度}に依存するとき、結合強度は
\(v_i\) と \(v_z\) の関係（同一軸か否か、間接結合の強度）に依存する。

§3.2 の \(f(Q)\) の世代依存パターンは、この結合強度が空間軸スケール
\(v_i\)
の非線形関数であることを示唆しており、正しい質量公式の判定条件となりうる。

{[}4{]} 定義9.4--9.5 の写像関数との対応を整理すると:

\begin{itemize}
\tightlist
\item
  \textbf{\(f_m^{\text{rest}}\)}（静止質量写像）: 本稿 §2--§3
  の考察対象。ヒッグス定常波との結合に由来する質量であり、フェルミオンに適用される。§2.1
  の \(C \cdot |v_i \cdot v_j|\) が最も単純な候補、§3.1 の
  \(f(Q, v_i) \cdot |v_i \cdot v_j|\) がQ依存性を含む候補である。§3.2
  により、\(f_m^{\text{rest}}\) は \(Q\)
  軸と空間軸の非自明な結合を含む。
\item
  \textbf{\(f_m^{\text{curv}}\)}（曲率エネルギー写像）: \(W^{\pm}, Z^0\)
  等のゲージボソンに適用される。{[}5{]} §8 の \(\delta k_H\)
  による質量修正は \(f_m^{\text{rest}}\)
  の範疇であるが、ゲージボソンの質量は主観空間の曲率半径 \(R_1\)
  に由来する加速度エネルギーであり（{[}5{]}
  補講4）、\(f_m^{\text{curv}}\) に対応する。\(f_m^{\text{curv}}\)
  の関数形は本稿の範囲外である。
\end{itemize}

\subsubsection{5.2
検証可能な構造的予測}\label{ux691cux8a3cux53efux80fdux306aux69cbux9020ux7684ux4e88ux6e2c}

質量公式の関数形に依存せず、本枠組みの構造から直接従う予測がある:

\textbf{予測 1}: 同一セクター（同一 \(Q\) 値、同一 \(t\)
符号）内の世代間質量比は、空間軸スケール比 \(v_x : v_y : v_z\)
のみで決まる。レプトンで \(v_x : v_y : v_z\)
を決定すれば、各クォークセクター内の質量比パターン（最重/最軽の比、中間/最軽の比）が同一になるかどうかが検証可能である。

\textbf{予測 2}:
ニュートリノのフレーバー間質量比は荷電レプトン質量比と特定の関係（§2.3
の観察2.2）を持つ。ニュートリノ質量の精密測定により検証可能である。

\subsubsection{5.3
本考察の限界}\label{ux672cux8003ux5bdfux306eux9650ux754c}

本稿は以下を確認した:

\begin{enumerate}
\def\labelenumi{\arabic{enumi}.}
\tightlist
\item
  軸スケール値 \(v_x, v_y, v_z, v_t\)
  による質量の記述は、パラメータ間に強い構造的拘束を与える
\item
  最も単純な質量公式（\(m \propto v_i \cdot v_j\), \(f(Q)\)
  定数）は定量的に不十分である（§3.2, §3.3）
\item
  Q軸と空間軸の結合関数 \(f(Q, v_i)\) の関数形の決定が残された課題である
\item
  \(v_x \neq v_y\)（\(m_\mu / m_e \approx 206.8\)）となる機構については、\(\mathrm{SO}(2)\)
  残留対称性のみからは説明されない。{[}4{]} §9.11.1 で「\(m_2 \neq m_1\)
  の差を本枠組み内で導出することはできない」とされた \(x, y\)
  軸間の非対称性の構造的起源は、本稿の範囲を超える
\item
  観察2.2
  のニュートリノフレーバー間質量比は、将来の精密測定により検証または反証される可能性がある。反証された場合、\(m \propto v_i \cdot v_j\)
  という基本仮定そのものの再検討が必要になる
\end{enumerate}

これらは本格的な質量分析の前提条件の整理であり、質量値の具体的な導出は本稿の範囲外である。

\begin{center}\rule{0.5\linewidth}{0.5pt}\end{center}

\subsection{参考文献}\label{ux53c2ux8003ux6587ux732e}

\begin{itemize}
\tightlist
\item
  {[}3{]} 木原 範昭,
  ``離散空間における中心投影の全球被覆と5次元背景空間の整数論的必然性,''
  Zenodo preprint, DOI: 10.5281/zenodo.19624957 (2026).
\item
  {[}4{]} 木原 範昭, ``6次元超直方体の集合構造とその組合せ論的性質,''
  Zenodo preprint, DOI: 10.5281/zenodo.19721125 (2026).
\item
  {[}5{]} 木原 範昭,
  ``形不変波の相互作用------軸方向変位転送・波束変形・因果の遡及的構成,''
  Zenodo preprint, DOI: 10.5281/zenodo.19721128 (2026).
\item
  {[}S1{]} 木原 範昭,
  ``符号付き面積の定式化------電荷構造の導出とスピン2配置の帰結,'' 補講1
  (2026).
\end{itemize}
\end{document}
